\subsection{6. Ressourcenbilanz und der jetzt konkrete nächste
Nachweis}\label{ressourcenbilanz-und-der-jetzt-konkrete-nuxe4chste-nachweis}

{\def\LTcaptype{none} % do not increment counter
\begin{longtable}[]{@{}
  >{\raggedright\arraybackslash}p{(\linewidth - 4\tabcolsep) * \real{0.3333}}
  >{\raggedright\arraybackslash}p{(\linewidth - 4\tabcolsep) * \real{0.3333}}
  >{\raggedright\arraybackslash}p{(\linewidth - 4\tabcolsep) * \real{0.3333}}@{}}
\toprule\noalign{}
\begin{minipage}[b]{\linewidth}\raggedright
Ressource
\end{minipage} & \begin{minipage}[b]{\linewidth}\raggedright
Herkunft im geprüften Vertrag
\end{minipage} & \begin{minipage}[b]{\linewidth}\raggedright
Was noch zu liefern wäre
\end{minipage} \\
\midrule\noalign{}
\endhead
\bottomrule\noalign{}
\endlastfoot
W, η, innere Darstellung, endliche Clock & Tatsächliche
Quellkonstruktion und erneuter exakter Replay & Physische Identifikation
des Trägers \\
Hnat mit Δ,g & Bestehender ausdrücklich festgelegter Modellvertrag &
Physische Auswahl und Kalibrierung \\
Zweizustandsblock β↔ρ & Neue vollständige Anwendung aller nativen Terme
& Ursprüngliche Vorbereitung von β oder ρ \\
Unabhängige X/Nb-Kontrollen & Gewährte Kontrollstufe, nicht aus fixes H
abgeleitet & Instrumente und Schaltvertrag \\
Q oder der gemeinsame Übergang V & Hier genau ausgeschriebene
zusätzliche Operation & Compilerwort oder Instrument, das diesen
Übergang wirklich erzeugt \\
G-triviale Referenz mit N\_R=0,4 & Minimaler zusätzlicher
Ladungsspeicher für diese Konstruktion & Träger, Präparation, Energie ER
und alle Ressourcenkosten \\
Anfang & Vakuum × scharfe Referenzladung vier & Herkunft dieses
Gesamtzustands \\
Ausgang & Gemeinsame Besetzungen und Referenzwechsel & Ausführbare
gemeinsame Auslesung, ohne stilles Ausspuren relevanter Information \\
\end{longtable}
}

Die Minimalität ist präzise begrenzt: Grad zwei ist der kleinste
zugelassene G-invariante paritätsgerade N-ändernde Systemgrad; zwei
verschiedene Referenzladungen sind die kleinste Dimension für den
einmaligen scharfen Ausgleich; die Dreizustandsmatrix ist der kleinste
von diesem konkreten Anfang durch die beiden nichtverschwindenden
Kopplungen erzeugte Raum. Das ist keine globale Minimierung aller
möglichen Universen oder Compiler.

\textbf{Nächster Herkunftsnachweis:} Ein vorhandener ursprünglicher
Quellbaustein muss einen getypten Adapter zu V einschließlich
Referenzzustand und gemeinsamer Auslesung liefern. Ein Vorschlag allein
aus den bereits N-neutralen Fockwörtern scheitert jetzt am bewiesenen
Erhaltungssatz und braucht keine weitere Wortsuche. Ein anderer getypter
Quellbaustein bleibt eine offene Möglichkeit; bloße E8-Grade oder
gleiche Matrizenabmessungen reichen als Adapter nicht aus.

\subsection{7. Reproduktion und
Beweisumfang}\label{reproduktion-und-beweisumfang}

Aus diesem Ordner:

\begin{verbatim}
python3 -B -W error verify.py --output results_normal.json
python3 -OO -B -W error verify.py --output results_optimized.json
cmp results_normal.json results_optimized.json
\end{verbatim}

Benötigt werden Python, NumPy, SciPy und SymPy. Der Replay liest nur die
eingefrorenen Eingaben und braucht weder das ursprüngliche Repository
noch Netzwerkzugriff. \codeword{freeze.py} ist nur das Herkunftswerkzeug
für die erste Archivierung und gehört nicht zum wissenschaftlichen
Replay.

Es gibt \textbf{261 explizite exakte Prüfbedingungen pro Modus},
zusätzlich zu den separat ausgewiesenen 17 erhaltenen nativen
Konstruktorguards. Die vier Clockguards sind in den 261 enthalten.
Normaler und \codeword{-OO}-Lauf erzeugen byteidentische JSON-Dateien.
Die Bosonidentitäten nutzen die exakte Weyl- Normalordnung; die nativen
Zustandsaktionen nutzen alle CAR-Vorzeichen und Bosonmultiplikitäten.
Keine Prüfung benutzt ein wegoptimierbares \codeword{assert}, und
Warnungen werden als Fehler behandelt.

Die Bedingungen prüfen konkrete algebraische Tatsachen und Quellenpins;
die allgemeinen Induktions- und Verfügbarkeitsargumente stehen oben. Die
Zahl 261 ist keine Zahl unabhängiger Theoreme oder physischer
Realisierungen. Kein neuer Grundzustandssatz für die Erweiterung, keine
native Raumzeit, keine vollständige E8-Operatordarstellung und kein
T1-T8- Abschluss wird behauptet. Fremde Quellen, v1.6.7, zentrale Ledger
und Webseiten wurden nicht verändert.

\begin{center}\rule{0.5\linewidth}{0.5pt}\end{center}

\section{Teilbericht B: vollständige innere
Schattenkarte}\label{teilbericht-b-vollstuxe4ndige-innere-schattenkarte}

\section{Gemeinsame Schattenkarte auf dem ursprünglichen
W}\label{gemeinsame-schattenkarte-auf-dem-urspruxfcnglichen-w}

Stand: 15. September 2026. Eigenständiger Folgeschritt zu v1.6.7.

\subsection{Ergebnis}\label{ergebnis-1}

Der ursprüngliche Tensor liefert jetzt eine vollständig berechnete
gemeinsame Schattenkarte für den inneren hellen N=2-Code. Einzelne
Besetzungen unterscheiden nur 24, sämtliche Einteilchenoperatoren nur
736 von insgesamt 3600 hermiteschen Operatorrichtungen. Paarbesetzungen
zusammen mit ausführbaren Drehungen aus G=Spin(10)×SU(4) liefern dagegen
eine vollständige Zustandsrekonstruktion. Eine konkrete endliche Familie
mit 3599 nichtkonstanten Erwartungswerten liegt bei. Diese letzte
Aussage ist \textbf{bedingt auf die Ausführbarkeit der zusätzlichen
Kontrollen und Messungen}. Aus ihrer algebraischen Existenz folgt noch
kein ursprüngliches Messinstrument samt Zustandsänderung.

Das Ergebnis betrifft denselben ursprünglichen W, keine ersatzweise
eingeführte Qubit-Analogie. Der W-Dateipin ist
\codeword{3f00a0892157a691e4ef26920c3702772c7bf76a250fad98ff04a079bb64b763}.
Die vier gelesenen v1.6.7-Kontextdateien sind ebenfalls unverändert
eingefroren.

\subsection{1. Ein gemeinsamer Träger, mehrere
Schatten}\label{ein-gemeinsamer-truxe4ger-mehrere-schatten}

Mit V=W†/√8 gilt V†V=I auf C\^{}60. Für eine logische Dichtematrix ρ
lautet der Schatten eines Fermionoperators O

\[\adjustbox{max width=\linewidth}{$\displaystyle 
  \Phi_{\mathcal O}(\rho)_O=\operatorname{tr}(\rho V^\dagger O V).
$}\]

Alle folgenden Messfamilien beziehen sich auf genau dieselbe ρ und
denselben Kodierer. Die Dimension der reellen Menge hermitescher
60×60-Matrizen ist 3600; bekannte Spur eins reduziert die freien
Zustandsparameter auf 3599. Alle angegebenen Familien enthalten die
Identität in ihrem linearen Raum, sodass ihre unsichtbaren Kerne bereits
spurfrei sind.

{\def\LTcaptype{none} % do not increment counter
\begin{longtable}[]{@{}
  >{\raggedright\arraybackslash}p{(\linewidth - 6\tabcolsep) * \real{0.2143}}
  >{\raggedleft\arraybackslash}p{(\linewidth - 6\tabcolsep) * \real{0.2857}}
  >{\raggedleft\arraybackslash}p{(\linewidth - 6\tabcolsep) * \real{0.2857}}
  >{\raggedright\arraybackslash}p{(\linewidth - 6\tabcolsep) * \real{0.2143}}@{}}
\toprule\noalign{}
\begin{minipage}[b]{\linewidth}\raggedright
Messfamilie
\end{minipage} & \begin{minipage}[b]{\linewidth}\raggedleft
Exakter Operatorrang
\end{minipage} & \begin{minipage}[b]{\linewidth}\raggedleft
Unsichtbare Zustandsdifferenzen
\end{minipage} & \begin{minipage}[b]{\linewidth}\raggedright
Verfügbarkeit
\end{minipage} \\
\midrule\noalign{}
\endhead
\bottomrule\noalign{}
\endlastfoot
Alle 64 Einzelbesetzungen n\_r & 24 & 3576 & Konkrete Operatoren;
ursprüngliche Messausführung offen \\
Alle hermiteschen Einteilchenbilineare f\_r†f\_s & 736 & 2864 &
Zusätzliche vollständige Bilinearfamilie \\
Alle G-gedrehten Einzelbesetzungen & 736 & 2864 & Bedingt auf
ausführbare G-Kontrollen \\
Alle 2016 ungedrehten Paarbesetzungen n\_i n\_j & 60 & 3540 & Gemeinsame
Besetzungsmessung erforderlich \\
Bilineare plus ungedrehte Paarbesetzungen & 772 & 2828 & Beide vorigen
Ressourcen zusammen \\
G-gedrehte Paarbesetzungen & 3600 & 0 & Gemeinsame Messung und
ausführbare G-Kontrollen \\
\end{longtable}
}

Der Rang zählt unabhängige lineare Erwartungswerte. Er zählt weder
Messapparaturen noch Messwiederholungen noch physische Raumrichtungen.

\subsection{2. Exakter Rangbeweis am
Tensor}\label{exakter-rangbeweis-am-tensor}

Die Datei enthält zehn symmetrische 16×16-Matrizen β\_k. Mit den sechs
antisymmetrischen 4×4-Farbmatrizen ε\_c gilt eintragsweise

\[\adjustbox{max width=\linewidth}{$\displaystyle 
 W_{(k,c),(p,a)(q,b)}=(\beta_k)_{pq}(\epsilon_c)_{ab}
 \quad\text{für }(p,a)<(q,b).
$}\]

Der Prüfer rekonstruiert hieraus den gesamten W und vergleicht alle
Einträge. Für sämtliche 4096 Bilineare vergleicht er zudem die direkte
CAR-Kontraktion mit der Faktorformel

\[\adjustbox{max width=\linewidth}{$\displaystyle 
 8V^\dagger f_{p,a}^\dagger f_{q,b}V
   =F_{pq}\otimes C_{ab},\qquad
 (F_{pq})_{kl}=\sum_t(\beta_k)_{pt}(\beta_l)_{qt},\quad
 (C_{ab})_{cd}=\sum_t(\epsilon_c)_{at}(\epsilon_d)_{bt}.
$}\]

In den aus den W-Faktoren abgeleiteten orthogonalen reellen Rahmen ist
der Spin-Bildraum genau \codeword{Skalar ⊕\allowbreak{} i·antisymmetrisch}, Dimension
1+45=46; der Farbbildraum hat Dimension 1+15=16. In der ursprünglichen
Gewichtsbasis prüft der Code die äquivalente ganzzahlige Identität

\[\adjustbox{max width=\linewidth}{$\displaystyle 
 n(A+\eta A^T\eta)=2\operatorname{tr}(A)I_n.
$}\]

Sie liefert jeweils die obere Ranggrenze. Ein nichtverschwindender
modularer Minor über dem Primkörper mit 1009 Elementen liefert die
passende untere Grenze. Damit sind die Ränge über Q und C exakt
bewiesen; es handelt sich nicht um eine toleranzabhängige numerische
Rangschätzung. Produktbildung ergibt 736.

Für die diagonalen Einzelbesetzungen bleiben nur Skalar plus Cartan:
(1+5)(1+3)=24. Die Adjungierten von so(10) und so(6) sind irreduzibel;
jede ursprüngliche Einzelbesetzung hat einen nichtverschwindenden
Skalar- und Cartananteil in beiden Faktoren. Ihr G-Orbit erzeugt deshalb
alle vier Summanden von (1⊕\allowbreak{}45)⊗(1⊕\allowbreak{}15), insgesamt 736, und keine
weiteren.

Der Eingabekern der gesamten Einteilchenkompression hat Dimension
\[4096-736=3360=(256-46)\cdot16.\] Er ist vom unsichtbaren \emph{Zustandskern}
mit Dimension 3600−736=2864 zu unterscheiden.

\subsubsection{Was genau fehlt?}\label{was-genau-fehlt}

Die gesamte Hermitesche Algebra zerfällt in den beiden reellen Faktoren
als

\[\adjustbox{max width=\linewidth}{$\displaystyle 
 (1\oplus i\Lambda^2\mathbb R^{10}\oplus\operatorname{Sym}^2_0\mathbb R^{10})
 \otimes
 (1\oplus i\Lambda^2\mathbb R^6\oplus\operatorname{Sym}^2_0\mathbb R^6).
$}\]

Die letzten symmetrischen, spurlosen Komponenten haben Dimension 54 und
20. Dem Bilinearschatten fehlen exakt die drei disjunkten Räume

\[\adjustbox{max width=\linewidth}{$\displaystyle 
 54\otimes(1\oplus15),\quad (1\oplus45)\otimes20,
 \quad54\otimes20,
$}\]

mit Dimension 864+920+1080=2864. Das ist eine konkrete Beschreibung des
vollständigen unsichtbaren Kerns, nicht bloß eine Dimensionsdifferenz.

Ein Verzicht auf alle imaginären Phasen wäre nochmals stärker: In
demselben reellen Gesamtrahmen haben reelle symmetrische Operatoren
Dimension 55·21+45·15=1830; imaginäre antisymmetrische Operatoren haben
Dimension 55·15+45·21=1770. Ein ausschließlich reeller Schatten könnte
diese 1770 Richtungen nicht sehen. Die native G-Familie enthält die
nötigen phasenempfindlichen Richtungen; ihre Ausführbarkeit wird damit
nicht bewiesen.

\subsection{3. Zwei orthogonale Zustände mit identischen eingeschränkten
Schatten}\label{zwei-orthogonale-zustuxe4nde-mit-identischen-eingeschruxe4nkten-schatten}

In der ursprünglichen logischen Gewichtsbasis, nullbasiert mit A=6k+c,
setze

\[\adjustbox{max width=\linewidth}{$\displaystyle 
 |\psi_\pm\rangle=\frac{|0\rangle\pm|30\rangle}{\sqrt2},
 \qquad \rho_\pm=|\psi_\pm\rangle\langle\psi_\pm|.
$}\]

Beide Matrizen sind exakt positiv, rein und zueinander orthogonal. Es
gilt für alle r,s und für alle i\textless j

\[\adjustbox{max width=\linewidth}{$\displaystyle 
 \operatorname{tr}[(\rho_+-\rho_-)V^\dagger f_r^\dagger f_sV]=0,
 \qquad
 \operatorname{tr}[(\rho_+-\rho_-)V^\dagger n_in_jV]=0.
$}\]

Insbesondere ändern beliebig viele G-gedrehte Einteilchenmessungen daran
nichts. Die Zustände unterscheiden sich in einer fehlenden
Spin-54-Richtung.

Der normierte logische Vektor

\[\adjustbox{max width=\linewidth}{$\displaystyle 
 |z\rangle=\tfrac12(|0\rangle+|30\rangle+i|6\rangle+i|36\rangle)
$}\]

liegt hingegen im G-Orbit eines ursprünglichen Kanalvektors. Die
zugehörige gedrehte Paarbesetzung komprimiert zu \textbar z⟩⟨z\textbar/8
und liefert die exakt verschiedenen Erwartungswerte 1/16 und 0. Dies ist
ein positiver Zustandszeuge für die Verbesserung durch die zusätzliche
Paarressource.

\subsection{4. Endliche, minimale Familie für vollständige
Rekonstruktion}\label{endliche-minimale-familie-fuxfcr-vollstuxe4ndige-rekonstruktion}

Jede unterstützte Paarspalte des W hat genau einen Eintrag±1. Daher gilt

\[\adjustbox{max width=\linewidth}{$\displaystyle 
 V^\dagger n_i n_jV=\tfrac18|A\rangle\langle A|
$}\]

für den zugehörigen Kanal A. Jeder der 60 Kanäle tritt auf. Ungedrehte
Paarbesetzungen erzeugen also exakt alle Diagonalmatrizen. Ihr Schnitt
mit dem Bilinearraum hat Dimension 24: Diagonalprojektion erhält die
oben angegebenen Faktor-Bildräume und hat in ihnen Rang 6 und 4. Damit
folgt der kombinierte Rang 736+60−24=772.

Der Prüfer leitet die reellen Rahmen aus den invarianten Faktormetriken
ab. Er prüft außerdem das W-Intertwining für sämtliche 45+15
Lie-Generatoren und ihre vollständigen Bilder so(10), so(6). G wirkt im
logischen Raum folglich als SO(10)×SO(6), jeweils mit seiner
Spin/SU(4)-Überlagerung.

Ein ursprünglicher Gewichtskanal ist in jedem Faktor ein zirkularer
Vektor (e\_i+i e\_j)/√2. Sein SO(n)-Orbit besteht aus (u+i v)/√2 mit
reellen, orthonormalen u,v. Für n=10 und n=6 erzeugt der Prüfer eine
rationale Liste solcher Projektoren:

\begin{enumerate}
\def\labelenumi{\arabic{enumi}.}
\tightlist
\item
  Alle (e\_i±i e\_j)/√2.
\item
  Für jedes i\textless j und ein drittes k den Vektor
  \codeword{(3e\_i+4e\_j+5i e\_k)/sqrt(50)}.
\end{enumerate}

Alle 50P haben ganzzahlige reelle und imaginäre Einträge. Der Code prüft
\codeword{P²=P}, \codeword{tr(P)=1} und die Isotropie
\codeword{tr(P Pᵀ)=0} exakt nach Skalierung. Die Isotropie und die Norm
kennzeichnen diese SO(n)-Gewichtsorbits; durch Orientierungswahl auf dem
orthogonalen Komplement existiert eine Drehung mit Determinante+1.
Beliebige U(10)- oder U(6)-Kontrollen werden nicht benutzt.

Modulare Zeilenauswahl liefert folgende konkrete unabhängige Basen:

\[\adjustbox{max width=\linewidth}{$\displaystyle 
 \mathcal B_{10}=\{I_{10},P_1,\ldots,P_{99}\},\qquad
 \mathcal B_6=\{I_6,Q_1,\ldots,Q_{35}\}.
$}\]

Die ausgewählten Vektoren und Indizes stehen vollständig im
Ergebnis-JSON. Ihre Produkte bilden eine Basis aller hermiteschen
60×60-Matrizen. Bei bekannter Spur bleibt eine minimale Familie von 3599
nichtkonstanten Zahlen:

\begin{itemize}
\tightlist
\item
  99·35=3465 gedrehte Paarprojektoren P\_a⊗Q\_b;
\item
  99 Spinwerte P\_a⊗I\_6, als Summe über die sechs Farbkanäle;
\item
  35 Farbwerte I\_10⊗Q\_b, als Summe über die zehn Spinkanäle.
\end{itemize}

Alle physisch angesetzten Paarerwartungswerte besitzen denselben Faktor
1/8; die Rekonstruktion kalibriert ihn ausdrücklich zurück. Da ein
offener Bereich um I/60 alle 3599 spurlosen hermiteschen Richtungen
enthält, kann keine feste Familie linearer Skalarerwartungswerte mit
weniger als 3599 Zahlen jeden Zustand identifizieren.
\textbf{Minimalität bezieht sich genau auf diese Zahlenzahl}, nicht auf
Geräteeinstellungen oder Messkosten.

